7. Limitations
7.1 Retrospective ERA5 Weather Inputs
All weather-informed model results in this study use same-hour ERA5 retrospective reanalysis weather input. ERA5 is a retrospective reanalysis product derived from historical data assimilation and is not an operational weather forecast. Same-hour ERA5 values are not available in real time; they represent conditions as reconstructed after the fact from a consistent global reanalysis system. All reported weather-informed model performance should therefore be interpreted as reflecting an idealized retrospective weather-information setting. These results do not establish what performance would be achievable under real-time operational conditions, where weather inputs would be subject to NWP forecast uncertainty.
7.2 No Operational NWP Forecast Comparison
This study does not quantify the degradation in model performance that would arise from replacing retrospective ERA5 inputs with real-time NWP forecast-weather inputs. Operational load forecasting requires weather inputs issued in advance of the target hour, which carry forecast uncertainty not present in the ERA5 retrospective setting used here. The extent to which point forecast accuracy and probabilistic coverage degrade under forecast-weather inputs is not evaluated and cannot be inferred from the results reported in Section 5. Future work could pair the modeling framework developed here with ensemble NWP forecast archives or operational forecast-weather inputs to quantify this source of uncertainty and assess performance in a setting that more closely reflects operational constraints.
7.3 Single-System and Single-Event Case Study
The results of this study are specific to PJM RTO aggregate hourly load and to the 2014 test year, with event-window analysis centered on the January 6–8, 2014 Polar Vortex case study. The January 6–8, 2014 Polar Vortex produced a near-annual-peak RTO load of 140,510.2 MW at Jan 7 18:00 EPT, equal to 99.18% of the 2014 annual peak of 141,677.9 MW recorded on Jun 17 17:00 EPT. Generalizability of the findings to other electricity systems, other ISOs or RTOs, other geographic regions, other years, or other categories of extreme events — including but not limited to heat waves, ice storms, and post-hurricane demand surges — is not established by this study and would require separate validation with appropriate data.
The evaluation covers a single out-of-sample test year. Whether the observed performance patterns, including the point forecast accuracy rankings and the event-window calibration failures, would replicate across multiple test years or under different training-period boundaries is not assessed. Multi-year evaluation across years with and without comparable cold-weather events would be required before broader conclusions about robustness could be drawn.
7.4 Quantile Crossing and Post-Hoc Rearrangement
The QR-GBT model produced quantile crossings in 5,786 of 8,760 full-year 2014 test hours, approximately 66% of all test hours. This high rate indicates that the seven independently fitted quantile models — each optimizing the pinball loss at a single quantile level without joint monotonicity constraints — do not form a coherent monotonic ordering across the majority of test hours. Post-hoc monotonic rearrangement was applied to restore a non-decreasing quantile order. This correction is a methodological limitation rather than a solution: rearrangement restores valid ordering but does not address the underlying cause of the crossing and does not guarantee well-calibrated prediction intervals.
The post-hoc character of the rearrangement correction means that the reported prediction intervals are an reflect both the original fitted quantile values and the rearrangement procedure. In applications where a coherent conditional distribution is required, architecturally enforced monotonicity — Architecturally constrained quantile models or distributional regression could enforce coherent predictive structure, while conformal prediction could be explored as a separate calibration framework — would be preferable to post-hoc correction.
7.5 Training-Distribution Support Not Directly Characterized
This study does not directly quantify whether the 2010–2013 training period contains cold-air events of intensity comparable to the January 2014 Polar Vortex. The observed degradation in point forecast accuracy and probabilistic coverage during the vortex window is consistent with distribution-shift stress, but a direct characterization of the training distribution in terms of cold-event frequency, temperature extremes, or load exceedance percentiles has not been conducted. Claims about the rarity or absence of comparable events in the training data should be treated as unverified. A future study could explicitly characterize training-distribution support by constructing multi-year weather and load percentile profiles, identifying historically comparable cold-air events in the training period, and assessing their relationship to observed model performance patterns.
7.6 Scope of Benchmark Comparison with PJM Day-Ahead
The PJM Day-Ahead forecast is included as an external operational benchmark and was produced using PJM's internal forecasting system and methods not controlled by this study. GBoost and QR-GBT use same-hour ERA5 retrospective reanalysis weather input, which gives these models access to weather information that was not available to PJM at the time of forecast issuance. The lower MAE values reported for GBoost relative to PJM Day-Ahead — both full-year (721 MW versus 1,937 MW) and during the vortex window (1,705 MW versus 3,148 MW) — therefore reflect this retrospective informational advantage and should not be interpreted as evidence that the proposed models would outperform PJM Day-Ahead under operationally equivalent conditions. Comparisons with the PJM Day-Ahead benchmark are retrospective in character and are included solely to contextualize the magnitude of modeled errors against those of an operational forecasting system during the same period.
7.7 Excluded Interpretability and External Data Claims
SHAP analysis. No SHAP-based feature attribution or feature importance analysis has been computed or validated for this study. No claims about the relative contributions of individual features to model predictions, nor any quantile-specific driver attributions, are made in this manuscript. Any such content from prior or legacy draft versions of this manuscript is excluded.
NOAA and NASA inputs. NOAA ISD station data are not used as model inputs and should not be treated as primary evidence in this manuscript. NASA AIRS imagery is not used, and no NASA weather data serve as inputs to any model evaluated here. Any legacy figures or references involving NASA AIRS imagery or NOAA ISD station-based weather are excluded from this study.
Deployment readiness. This is a retrospective research study. The models evaluated here have not been tested under real-time operational conditions, do not use operational forecast-weather inputs, and have not been validated against operational reliability requirements. No claims of deployment readiness, real-time operational superiority, or integration readiness with ISO or RTO systems are made or implied.
Implementation reproducibility. Exact hyperparameter configurations and software package versions for GBoost and QR-GBT should be documented in final reproducibility materials and confirmed against the analysis codebase before submission. If any metric implementation — including Winkler score or pinball loss variant — differs from the prose definitions provided in Section 4, the discrepancy should be reconciled in the final revision.